% app_d2_eps4.tex -- d = 2 appendix, part 1 of 3 (HT2dPert, 2026-10-16): operators and schemes at low order, solver
% checks, exact C4(L), the mechanism of the ln L divergence class by class, and the complete window at eps^4.
% Merged from the method write-up of package I (public edition of 2026-10-14), which these appendices
% replace; superseded material is not carried over, and what Sec. d2:sec states is not repeated.
% Notation of Sec. d2:sec and Table tab:notation (level l, operator \hat l, sub-cutoff lambda, e_1 = E_1 - E_0).
% Numbers quoted in Sec. d2:sec use the macros of numbers_d2.tex (replay rows in repro_d2.tex); every other number is
% read from, or (where marked) evaluated from, the records named in the subsection (all shipped in package I).
\section{Two dimensions: operators, exact fourth order and the complete window}\label{htp:d2:app:eps4}

Appendices~\ref{htp:d2:app:eps4} to~\ref{htp:d2:app:legs} give the derivations and the supporting data behind
Section~\ref{htp:d2:sec}. This appendix covers the operators and schemes at low order, the solver checks, the exact $C_4(L)$,
the cause of its $\ln L$ divergence class by class, and the complete window at $\eps^4$. Appendix~\ref{htp:d2:app:eps5}
treats order $\eps^5$: the window at all orders, the double counting that rule $W_c$ removes, and the tail of $C_5$.
Appendix~\ref{htp:d2:app:legs} treats the descendant spacings and the external-line windows $W_{cs}$ and $W_{cs4\chi}$, and
ends with the status of each window component. Paths are relative to the package
\code{code/\allowbreak{}lri\_2d\_coeff\_independent\_check/} (package I of Section~\ref{htp:sec:repro}), and each subsection names its
records. Numbers that Section~\ref{htp:d2:sec} also quotes have reproduction rows in Section~\ref{htp:sec:repro}. The other numbers
of these appendices are read from the records named in each subsection, or evaluated from them where the text says
so; the reproduction scripts do not recompute them. ``Exact'' data
are rationals, or float64 evaluations of exact finite sums with a measured rounding error of at most
$1.4\times10^{-11}$. Limits $L\to\infty$ are fits; their uncertainties are spreads over tail models and fit windows, not
standard deviations.

\subsection{Operators and schemes at low order}\label{htp:d2:app:ops}
{\small\raggedright\noindent Records: \code{SPEC.md}, \code{rs\_exact/\allowbreak{}README.md}, \code{zeromode\_exact/\allowbreak{}README.md},
\code{cauchy/\allowbreak{}README.md}.\par}

\paragraph{Factorized vertices.} With $Z=a_0+a_0^\dagger$, $a_0=A_{00}$, and $\chi=\phi-Z$, normal ordering factorizes:
\begin{equation}
 V=\sum_{m=0}^4\frac1{12}\binom4m{:}Z^{4-m}{:}\otimes\langle{:}\chi^m{:}\rangle_\theta,\qquad
 M=\sum_{m=0}^2\binom2m{:}Z^{2-m}{:}\otimes\langle{:}\chi^m{:}\rangle_\theta .
 \label{htp:d2:eq:factorize}
\end{equation}
At $\eps=0$ the weights factorize holomorphically, $w_{ab}=v(a)v(b)$ with $v(a)=(\frac12)_a/a!$. The squared
$V$-vertex coefficient times the Wick count of a bundle of $m$ $\chi$ lines is
$\gamma_m=m!\,\big[\binom4m/12\big]^2$: $\gamma_1=\frac19$, $\gamma_2=\frac12$, $\gamma_3=\frac23$, $\gamma_4=\frac16$.

\paragraph{Space and particle cap.} The truncated space holds the occupation states with spin zero, even under
$(a,b)\to(b,a)$, and $\ell\le L$; spin zero forces $\ell$ even, and the field parity $(-1)^P$ is conserved. Truncated
operators are the projections of the full ones; nothing else is truncated. The level cutoff does not bound $P$, because
zero-mode quanta cost no level. A path argument shows that a contribution to the order-$\eps^n$ energy of $|p\rangle$
reaches at most $P=p+4\lfloor n/2\rfloor$, which gives the cap $P\le p+16$ through $\eps^8$. It was also tested: caps
$p+18$ and $p+20$ reproduce the $p+16$ rationals identically, while $p+14$ changes $C_8$ and nothing below. The Cauchy
code confirms it with absolute caps: $P\le21$ and $P\le23$ reproduce $P\le19$, and $P\le17$, 18 change only $C_8$.

\paragraph{References and the projection $\Pi$.} Each reference $|p\rangle$ is non-degenerate in its sector, so $E_p(\eps)$
is unique. Because $w_{00}=1$ and every window has vanishing diagonal elements on level-0 states, $\langle p|M|p\rangle=2p$ and
$\langle p|V|p\rangle=\frac12p(p-1)$ at every $\eps$. With $\mathcal E_p^{(n)}$ the order-$n$ energy without the $r_n$,
$g_n$ terms, Eq.~\eqref{htp:eq:Pi} reads $C_n=\Pi[\mathcal E^{(n)}]$ with
\begin{equation}
 \Pi f=f(2)-f(1)-\tfrac13\big[f(3)-f(0)\big]=-\tfrac13\Delta^3f(0),\qquad \Pi[p^3]=-2,\qquad \Pi[p^4]=-12 ,
 \label{htp:d2:eq:Pi}
\end{equation}
and $\Pi$ annihilates every polynomial of degree at most two in $p$. Only the part of an energy shift that is at least
cubic in $p$ reaches $C_n$.

\paragraph{Where the schemes differ.} The diagonal windows (consistent, literal, zero-mode) are of order $\eps^2$ and
vanish on $|p\rangle$, so they first enter through $\langle p|H_1G_pD_2G_pH_1|p\rangle$ at $\eps^4$, with $H_1=F_1+g_1V$ and the resolvent
$G_p=Q_p/(\frac p2-H_0)$, $Q_p$ the projection off $|p\rangle$. Raw and consistent therefore differ first at $C_4$. Raw
and consistent vacua agree through $\eps^3$, so the consistent and literal windows differ first at $C_6$; the exact data
confirm both statements for every $L>0$. At $\eps^4$ only the $\eps^2$ parts of the windows matter:
\begin{equation}
 D^{\rm cons}_2(\ell)=E_v^{(2)}(L)-E_v^{(2)}(L-\ell),\qquad
 E_v^{(2)}(\lambda)=g_1^2\langle0|VG_0V|0\rangle_\lambda=-\frac{(\lambda+2)(\lambda+4)}{864},
 \label{htp:d2:eq:Ev2}
\end{equation}
where $\langle\cdot\rangle_\lambda$ restricts the intermediate states to level $\le\lambda$, and the zero-mode window is
$D_2^{\rm zm}(z,\ell)=h^{(2)}_{zz}(L)-h^{(2)}_{zz}(L-\ell)$, the diagonal of Eq.~\eqref{htp:d2:eq:Dzm}. The $r_2$ and $g_2$ terms of
the sub-cutoff energies cancel in these differences, and for $z=0$ the two windows coincide.

\paragraph{Degenerate partners.} For $z\ge5$ the zero-mode channel $|z\rangle$ has degenerate partners with level
$\ell=2,4,\dots\le\frac{z-1}2$ and $P=z-2\ell$. First order lifts the degeneracy (smallest separation $\frac73$, at $z=5$,
$\ell=2$), so the branch through $|z\rangle$ is unique; the references are never degenerate. The window objects project
the partners out (Appendix~\ref{htp:d2:app:Wall}).

\subsection{Solver checks in detail}\label{htp:d2:app:checks}
{\small\raggedright\noindent Records: \code{compare\_with\_package.txt}, \code{compare\_three\_way.txt}, \code{rs\_exact/\allowbreak{}README.md},
\code{cauchy/\allowbreak{}README.md}, the production records \code{../\allowbreak{}../\allowbreak{}data/\allowbreak{}HT\_2d/\allowbreak{}prod/}, \code{continuum/\allowbreak{}exact\_C2\_C3.csv}, \code{continuum/\allowbreak{}asymptotics\_C2\_C3.txt}; code
\code{continuum/\allowbreak{}exact\_low\_orders.py}.\par}

\paragraph{The three computations.} Section~\ref{htp:d2:solvers} and the independent-checks paragraph of
Section~\ref{htp:sec:checks} give their scope and agreement. The production solver propagates Taylor jets of the wave
functions in float64, with scalar reductions in extended long double; its consistent run at $L=34$ took 2893\,s on 64
threads with a peak memory of 463\,GiB. The exact solver works in a similarity frame in which every matrix element is
rational; its modular code uses 72 primes with Chinese remaindering, and every rational reconstruction is verified on
eight held-out primes. The Cauchy solver evaluates $\DE$ at $\eps_k=\rho\,e^{2\pi\mathrm ik/K}$ with $\rho=\frac1{64}$,
$K=128$ and 256-bit arithmetic, and $C_n=K^{-1}\sum_k\DE(\eps_k)\eps_k^{-n}$. It solves the consistent vacua exactly. On states with
$\ell\ge2$ the window is $E_v(\lambda)-E_v(\lambda-\ell)$, so the self-consistent condition $(H_\lambda+D_\lambda)\psi=E_v(\lambda)\psi$
becomes the linear generalized eigenproblem $A_\lambda\psi=E_v(\lambda)\,\mathcal P_0\psi$, with
$A_\lambda=H_\lambda-E_v(\lambda-\hat\ell)(1-\mathcal P_0)$, $\mathcal P_0$ the projection on level-0 states, and the
sub-cutoff vacua $E_v(\lambda-\ell)$ known from smaller $\lambda$.

\paragraph{Internal checks of the exact codes.}
\begin{itemize}
\item $C_0=1$, $C_1=-\frac16$, $r_1=0$, $g_1=\frac13$ for every $L$ and scheme, and the tuning conditions hold exactly at
 every order. The fraction and modular codes agree identically for every number at $L\le12$.
\item A third, Hermitian construction of $H$ was diagonalized at 110 digits at $\eps=10^{-3}$, $10^{-5}$, $10^{-8}$, with the
 exact series for $r$, $g$ and the vacua inserted ($L\le6$, every scheme and channel). The residual divided by $\eps^9$
 tends to a constant, which confirms every coefficient through $\eps^8$ to about $10^{-8}$.
\item The Cauchy results do not change with the circle radius, the number of points, the precision, the continuation path
 or the particle cap; the quoted errors (about $10^{-45}$ on $C_8$) are set by the most aliased variant. All 882 rational
 identifications have denominators of the form $2^a3^b5^c7^d$.
\end{itemize}

\paragraph{$C_2$ and $C_3$.} The closed form~\eqref{htp:d2:eq:C2} gives $C_2(0)=\frac{17}{81}$, $C_2(2)=\frac{19}{72}$,
$C_2(4)=\frac{179}{640}$ and $C_2(6)=\frac{833719}{2903040}$, equal to the exact rationals (evaluated against
\code{rs\_exact/\allowbreak{}results/\allowbreak{}raw\_L\{0,2,4,6\}.json}). Its approach to the limit is
$C_2(L)-C_2=-\frac{40}{81\pi}L^{-1}+0.235785\,L^{-2}+\dots$. An independent diagrammatic code, which contracts the factorized
vertices~\eqref{htp:d2:eq:factorize} with the level cutoff on every intermediate state, gives $C_3(L)$ in float64 to
$L=\NTwoCthreeLmax$, where $C_3(600)=-0.865794525554$. Its $\ln L/L$ coefficient is $\NTwoCthreeTail$, which is
$4.4999\times\frac{40}{81\pi}$; the factor $\frac92$ is identified numerically from this fit. It is not the one-loop $\phi^6$
value $\omega_1=\frac72$ of Eq.~\eqref{htp:eq:omega}, and which operator sets the $d=2$ tail is not derived. With the tower
$1$, $\ln L/L$, $1/L$, $\ln^kL/L^2$ ($k\le2$), $\ln L/L^3$, $1/L^3$ the limit is $-0.8748500372$, $-0.8748500106$ and
$-0.8748500065$ for $L_{\min}=32$, 64 and 128, which gives $C_3=\HTpCthreeTwo$.

\subsection{Exact \texorpdfstring{$C_4(L)$}{C4(L)} to \texorpdfstring{$L=2000$}{L=2000}}\label{htp:d2:app:c4exact}
{\small\raggedright\noindent Records: \code{c4\_exact/\allowbreak{}README.md}, \code{c4\_exact/\allowbreak{}C4\_by\_L.csv},
\code{c4\_exact/\allowbreak{}results/\allowbreak{}\{c4\_L1000,\allowbreak{}c4\_L2000,\allowbreak{}windows\_L20000\}.npz}, \code{c4\_exact/\allowbreak{}fits/\allowbreak{}fit\_summary.txt},
\code{c4\_exact/\allowbreak{}fits/\allowbreak{}blind\_predictions\_L1200\_2000.json}, \code{c4\_exact/\allowbreak{}validate\_\{rs,\allowbreak{}pkg,\allowbreak{}corrected\}.txt},
\code{zeromode\_exact/\allowbreak{}README.md}, \code{continuum/\allowbreak{}\{fits\_summary,\allowbreak{}zm\_corrected\_fits\}.txt}; code \code{c4\_exact/\allowbreak{}\{c4lib,\allowbreak{}c4\_compute,\allowbreak{}c4\_run,\allowbreak{}windows,\allowbreak{}fitlib,\allowbreak{}fits\_c4\}.py}.\par}

At $L\le30$ the production data could not separate a slow power tail from a logarithm. Fits of the consistent $C_4$ with
towers that contain $\ln L/L$ landed 2.4 to 5\% below the continuum value, and the same fits with the zero-mode window added
gave 4.90 to 5.04 (\code{continuum/\allowbreak{}zm\_corrected\_fits.txt}). The exact $C_4(L)$ settles the question.

\paragraph{Fourth-order perturbation theory.} Write $H=H_0+\sum_{k\ge1}\eps^kH_k$ with
\begin{equation}
 H_1=F_1+g_1V,\qquad H_2=r_2M+g_2V+g_1\dot V+D_2,\qquad H_3=r_3M+g_3V+g_2\dot V+\tfrac12g_1\ddot V+r_2\dot M+D_3 ,
\end{equation}
where $\dot V=\partial_\eps V$ and $\dot M=\partial_\eps M$ act through $w_{ab}(\eps)$. With $G=G_p$, $\langle\cdot\rangle=\langle
p|\cdot|p\rangle$ and $E^{(1)}=\langle H_1\rangle$,
\begin{align}
 E^{(4)}={}&\langle H_4\rangle+2\langle H_1GH_3\rangle+\langle H_2GH_2\rangle+2\langle H_1G(H_1-E^{(1)})GH_2\rangle
 +\langle H_1G(H_2-\langle H_2\rangle)GH_1\rangle\nonumber\\
 &+\langle H_1G(H_1-E^{(1)})G(H_1-E^{(1)})GH_1\rangle-\langle H_1GH_1\rangle\langle H_1G^2H_1\rangle .
\end{align}
A diagonal window enters $C_4$ only through $g_1^2\langle VGD_2GV\rangle$; the mass window $D^M$ enters through more
terms (Appendix~\ref{htp:d2:app:window}). The order-3 couplings $r_3$, $g_3$ need the complete
$\mathcal E^{(3)}_p$ for $p=0,1,3$, not only its part that survives $\Pi$.

\paragraph{Diagram evaluation.} With Eq.~\eqref{htp:d2:eq:factorize}, a $K$-vertex chain is a sum over $\chi$-line patterns
$n_{ij}$ (lines between vertices $i<j$) of a Wick count, the vertex coefficients, an exact rational zero-mode path sum and
a $\chi$ level sum. A time slice $s$ with $n_s$ $\chi$ lines and $z_s$ zero-mode quanta has the energy denominator
$\frac12(z_s-p)+\frac12n_s+(\text{level of the slice})$. The $\chi$ sums use the holomorphic factorization: with
$v_k(y)=(1-y)^{-k/2}$, $\sum_{ab}w_{ab}y^a\bar y^b=v_1(y)v_1(\bar y)$, and a bundle of $n$ $\chi$ lines has the generating
function $[v_1(y)v_1(\bar y)-1]^n=\sum_k\binom nk(-1)^{n-k}v_k(y)v_k(\bar y)$. Angular averaging at every vertex forces spin
zero in every slice, so the holomorphic and antiholomorphic levels of a slice are equal, $j_s$, and its level is $2j_s$.
The density of a pattern is then a sum of squares of one-variable coefficients, with no angular-momentum sum left.
$\eps$-derivative vertices differentiate the line weights.

\paragraph{One run for all cutoffs.} The densities do not depend on $L$; the cutoff only imposes $\max_s2j_s\le L$.
Accumulating by $\max_sj_s$ and summing cumulatively gives every even $L\le L_{\max}$ in one run. The window terms
depend on $L$ explicitly and are evaluated separately.

\paragraph{Pruning and symmetry.} The zero-mode amplitude of a four-$V$ chain with $n$ $\chi$ lines is a polynomial of
degree $8-n$ in $p$. Since $\Pi$ kills degree $\le2$, only $n\le5$ reaches $C_4$: 384 of the 641 four-vertex patterns.
Time reversal maps them onto each other in 168 pairs with identical contributions, so 216 patterns are computed. The full set reproduces every
individual $E_p^{(4)}$ of the exact solver.

\paragraph{Window shifts without cancellation.} Evaluated naively, the window $h^{(2)}_{zz}(L)-h^{(2)}_{zz}(L-\ell)$ is built from
terms of size $h^{(2)}_{zz}\sim L^2$, while the part that survives $\Pi$ is of order $1/L$: a relative cancellation of about
$L^3$. The script \code{windows.py} applies $\Pi$ to the
rational zero-mode polynomials first and only then multiplies by single-signed level sums. It runs to $L=20000$ and agrees
with the full computation to $6\times10^{-13}$ for $L\le1000$.

\paragraph{Cost and precision.} On one 40-core node the run to $L_{\max}=600$ takes 51\,s, to 1000 four minutes and to
2000 about 77 minutes (two Slurm jobs on the symmetry cluster, 59 minutes to the time limit and 17 minutes from the
checkpoint; the package README's 80 to 90 minutes is the estimate made during the run). Everything is float64. Overlapping
runs agree to $6.7\times10^{-12}$ (600 against 1000, from the README) and $1.4\times10^{-11}$ (1000 against 2000, evaluated
from \code{c4\_L1000.npz} and \code{c4\_L2000.npz}), and the cancellation in the assembly is about $4\times10^3$ at $L=1000$.

\begin{table}[ht]
\centering\footnotesize\renewcommand{\arraystretch}{1.15}
\begin{tabular}{@{}p{8.6cm}p{2.2cm}p{3.7cm}@{}}
\toprule
check & range & worst deviation \\\midrule
$r_2$, $g_2$, $r_3$, $g_3$, $C_2$, $C_3$ against the exact rationals & $L=0$ to 18 & $3.4\times10^{-15}$ relative ($C_3$: $1.3\times10^{-14}$) \\
raw and consistent $C_4$ against the exact rationals & $L=0$ to 18 & $1.1\times10^{-13}$ relative \\
every $E_p^{(4)}-2pr_4-\frac12p(p-1)g_4$, $p=0$ to 3, raw, all 641 patterns & $L=0$ to 18 & $2.1\times10^{-14}$ relative \\
$E_v^{(2)}(\lambda)$ of Eq.~\eqref{htp:d2:eq:Ev2} & $\lambda\le18$ & $2.2\times10^{-16}$; exact for $\lambda\le12$ \\
$C_2$, $C_3$, raw and consistent $C_4$, $r_3$, $g_3$ against the production solver & $L=2$ to 34 & $1.1\times10^{-13}$ relative \\
zero-mode $C_4$ against brute-force fraction arithmetic & $L=0$ to 12 & $1.1\times10^{-13}$ relative \\
zero-mode $C_4$ against the exact zero-mode solver & $L=0$ to 18 & $1.1\times10^{-13}$ relative \\
pruned ($n\le5$) against the full pattern set & $L\le18$ & $10^{-13}$ \\
time-reversal pairs & 168 pairs, $L=40$ & $5.5\times10^{-16}$ \\
\bottomrule
\end{tabular}
\caption{Validation of the large-$L$ $C_4$ code (\code{c4\_exact/\allowbreak{}README.md}, \code{c4\_exact/\allowbreak{}validate\_*.txt}).}
\label{htp:d2:tab:c4valid}
\end{table}

\paragraph{Validation.} Table~\ref{htp:d2:tab:c4valid} lists the checks. The zero-mode scheme itself was solved exactly through
$\eps^8$ for $L=0$ to 18 (\code{zeromode\_exact/}), by exact perturbation theory with the degenerate partners treated by
a projected equation. Two implementations (modular and fraction) are identical for $L\le10$. A non-perturbative Cauchy
check covers every coefficient and every sub-cutoff energy $E_z^{(k)}(\lambda)$ at $L\le6$, and closure tests show that this scheme
needs only the zero-mode channels $z\le14$ through $\eps^8$. The zero-mode $C_4$ minus the consistent $C_4$ equals an exact
closed form in $L$ (\code{zeromode\_exact/\allowbreak{}closed\_form\_dC4.py}) at every $L$; for example
$C_4^{\rm zm}(2)=\frac{6786263}{3732480}$.

\paragraph{Exact window shifts.} Fits of the cancellation-free shifts to $L=20000$ (residuals $\le5\times10^{-14}$) give
\begin{align}
 \Delta C_4^{\rm cons}(L)&\equiv C_4^{\rm cons}-C_4^{\rm raw}=\kappa L-\frac5{8748\pi}\ln L+0.000446687+o(1),\nonumber\\
 \Delta C_4^{\rm zm}(L)&\equiv C_4^{\rm zm}-C_4^{\rm cons}=\Big(\frac4{81\pi}+\frac5{4374\pi}\Big)\ln L+0.0187584+o(1),
 \label{htp:d2:eq:dC4shifts}
\end{align}
with $\kappa=\NTwoKappa$; the analytic evaluation in \code{continuum/\allowbreak{}asymptotics\_C2\_C3.txt} gives
$\kappa=2.20729299645\times10^{-4}$. The $\ln L$ coefficients are those of Table~\ref{htp:d2:tab:logs}. Split by the number of
outer $\chi$ lines, $\Delta C_4^{\rm zm}$ has logarithms only for two lines ($\frac4{81\pi}$) and four lines
($3.6387\times10^{-4}$). Evaluated from \code{C4\_by\_L.csv}, it has a minimum of $0.0982$ near $L=36$ and then grows logarithmically: $0.1028$ at $L=100$,
$0.1313$ at 1000 and $0.1418$ at 2000.

\paragraph{Fits.} The tail models (\code{fitlib.py}) are towers $T_{jk}$: a constant, $\ln^iL/L$ for $i\le j$ and
$\ln^iL/L^2$ for $i\le k$; $T_{jkm}$ adds $\ln^iL/L^3$ for $i\le m$. A model may add a free $\ln L$ and single terms such
as $\ln^3L/L$. The table of Section~\ref{htp:d2:c4} uses $T_{22}$ and $T_{22}+\ln^3L/L$, each with a free $\ln L$, the $\ln^2L/L$
coefficient free or fixed to $-5/\pi$ (end of Appendix~\ref{htp:d2:app:window}), and $L_{\min}\in\{200,300,400\}$, $L\le2000$: twelve fits, with the median and half
the full spread quoted. At $L=2000$ the data are $C_4=4.29126$ (raw), $4.73178$ (consistent) and $4.87359$ (zero-mode). The
raw fit uses $C_4^{\rm raw}=C_4^{\rm cons}-\Delta C_4^{\rm cons}$ with the exact shift. The difference of the two fitted
$\beta$, $0.016082$, equals the $\ln L$ coefficient of $\Delta C_4^{\rm zm}$. Subtracting the fitted $1/L$ and $1/L^2$
terms from $C_4(L)$ leaves straight lines in $\ln L$: $5.0026-0.02656\ln L$ (consistent) and $5.0213-0.01048\ln L$
(zero-mode). Without a $\ln L$ term every tower drifts with
the fit window, and the limits, 4.881 to 4.896 (zero-mode) and 4.648 to 4.684 (consistent), are far from $5.0489$ and from
each other. The fitted $\beta$ of the consistent scheme alone would not exclude the nearby simple value $-1/(12\pi)=-0.0265258$; the class
decomposition of Appendix~\ref{htp:d2:app:mech} fixes it at $-\frac{365}{4374\pi}$.

\paragraph{Out-of-sample and blind tests.} Fitted on $[200,600]$ and used to predict $(600,2000]$, the zero-mode model
$T_{22}$ with $\ln L$ misses by at most $9.0\times10^{-8}$, and models without $\ln L$ miss by $5.0\times10^{-6}$ to
$2.5\times10^{-5}$. Fitted on $[200,1000]$, the models with $\ln L$ miss by $8.5\times10^{-9}$ ($T_{22}$) and
$5.5\times10^{-8}$ ($T_{232}$), those without by $4.3\times10^{-7}$ to $6.7\times10^{-6}$. Before the $L=2000$ run finished,
predictions for $L=1200$ to 2000 were recorded from fits on $L\le1000$. At $L=2000$, prediction minus data is
$-8.6\times10^{-9}$ ($T_{22}$) and $-6.4\times10^{-8}$ ($T_{232}$) for the models with $\ln L$ in the zero-mode and
consistent schemes ($-7.5\times10^{-9}$ and $-8.3\times10^{-8}$ raw). For the models without $\ln L$ it is $+4.3\times10^{-7}$
to $+1.8\times10^{-6}$ (zero-mode), $+1.2\times10^{-6}$ to $+4.6\times10^{-6}$ (consistent) and $-4.1\times10^{-7}$ to
$-8.6\times10^{-7}$ (raw). The data are accurate to about $10^{-11}$, so the blind test favors the logarithm, by factors of 5 to
500 (the ratios of the deviations above).

\subsection{Mechanism and the \texorpdfstring{$\ln L$}{ln L} coefficients by class}\label{htp:d2:app:mech}
{\small\raggedright\noindent Records: \code{residual\_log/\allowbreak{}results/\allowbreak{}log\_coeffs.txt}; code \code{residual\_log/\allowbreak{}\{log\_coeffs,\allowbreak{}offdiag\_window,\allowbreak{}decompose\}.py}.\par}

\paragraph{The budget defect.} The cutoff constrains the total level of each intermediate state. View a fourth-order
process as an outer structure with a sub-bubble nested in one of its intermediate states. The sub-bubble may use only the
level $L-\ell_{\rm other}$, where $\ell_{\rm other}$ is the level of everything else in that time slice, while the
counterterm that cancels its divergence was tuned on level-0 references, where it has the full $L$. The mismatch is the
density of the divergence at the cutoff times the level taken by the other quanta. Table~\ref{htp:d2:tab:budget} lists the
three divergent sub-bubbles of the $\phi^4$ vertex.

\begin{table}[ht]
\centering\footnotesize\renewcommand{\arraystretch}{1.15}
\begin{tabular}{@{}p{3.1cm}p{4.3cm}p{3.7cm}p{4.1cm}@{}}
\toprule
sub-bubble & divergence & deficit for spectators of level $\ell$ & restored by \\\midrule
vacuum (four $\chi$ lines) & $E_v^{(2)}(\lambda)=-\frac{(\lambda+2)(\lambda+4)}{864}$, quadratic & ${\sim}\ell L$ & consistent window \\
sunset (three $\chi$ lines; mass) & $r_2(\lambda)=\frac{2\lambda}{27\pi}+\dots$, linear & $\frac2{27\pi}\ell$ per mass insertion; does not decay & in part by the diagonal zero-mode window; fully by $D^{\rm zm}$ and $D^M$ \\
coupling (two $\chi$ lines) & $g_2\sim\ln\lambda$ & $\ln\frac L{L-\ell}\sim\frac\ell L$ & not needed for soft spectators; finite effects from hard ones \\
\bottomrule
\end{tabular}
\caption{Divergent sub-bubbles and their level-budget deficits in $d=2$ ($\Ds=\frac12$ at $\eps=0$).}
\label{htp:d2:tab:budget}
\end{table}

The sunset slope is $dr_2/dL=\frac2{27\pi}=0.0235785$ (the exact $[r_2(2000)-r_2(1998)]/2$ is $\NTwoRtwoSlope$). Because the
divergence is linear, the deficit is proportional to the spectator level and independent of $L$, and it hits every
insertion of the mass operator $M$, not only the diagonal of its zero-mode part.

\paragraph{Linked clusters with zero-mode transfer.} Take two components that share only the zero mode: an outer bubble $O$
with $m$ $\chi$ lines and an inner bubble $B$ with $m'$. Let $E_O$, $E_B$ be their slice energies and $q=\frac12\Delta z_B$ the
zero-mode transfer, with $\Delta z_B$ the change of $z$ caused by $B$. The four vertices
have six time orderings: two sequential ones ($OOBB$, $BBOO$), in which no slice holds lines of both components, and four
overlapping ones ($OBBO$, $BOOB$, $OBOB$, $BOBO$), in which some slice holds both. Since $H_0=\ell+\frac P2$ is additive, the
six orderings and the unlinked terms of perturbation theory cancel in the continuum. At leading degree in $p$ the zero-mode
amplitude does not depend on the ordering, and the cancellation reduces to the identity
\begin{equation}
 \frac1{E_OE_B(E_O+q)}+\frac1{E_OE_B(E_B-q)}-\frac1{q\,E_O(E_B-q)}+\frac1{q\,E_B(E_O+q)}=0
\end{equation}
for every $q$, including the limit $q\to0$; its first two terms come from the four overlapping orderings, the last two from
the sequential orderings and the unlinked terms. With outer level $2a$ and inner level $2b$, the overlapping orderings need
$a+b\le\frac L2$, and the others do not. The truncation therefore errs by
\begin{equation}
 -\sum_{a,b\le L/2<a+b}\Big[\frac1{E_OE_B(E_O+q)}+\frac1{E_OE_B(E_B-q)}\Big]\times(\text{zero-mode amplitude}).
\end{equation}
This is the nested insertion, into every outer state of level $\ell$, of the window part of the second-order zero-mode
effective interaction,
\begin{equation}
 D_\lambda(z\to z')=g_1^2\sum_{k:\,\lambda-\ell<\ell_k\le\lambda}\frac{\langle z'|V|k\rangle\langle k|V|z\rangle}{E^0_z-E^0_k}
 =h^{(2)}_{z'z}(\lambda)-h^{(2)}_{z'z}(\lambda-\ell),
 \label{htp:d2:eq:Dwindow}
\end{equation}
in its ket-referenced (Bloch) form and including $z'\neq z$: the window $D^{\rm zm}$ of Eq.~\eqref{htp:d2:eq:Dzm}. The diagonal
zero-mode scheme keeps only $q=0$. At subleading degree in $p$ the orderings differ by commutator terms. These are confined to
the window region and are $\mathcal O(\ln L/L)$ for the numbers $(m,m')=(2,2)$, $(1,3)$ and $(3,1)$ of outer and inner $\chi$
lines.

\paragraph{The ratio $\frac23$.} Take the outer two-$\chi$ bubble, with ${:}Z^2{:}$ at each outer vertex, and the inner
zero-mode sunset, with ${:}Z{:}$ at each inner vertex. Per unit amplitude the window of the nested sunset at outer level $2a$
is $-\frac2\pi a$, and the outer weight (pattern density over the squared denominator) tends to $\frac1{4a^2}$. The window sum
is therefore
\begin{equation}
 \sum_{a\le L/2}\frac1{4a^2}\Big(-\frac{2a}\pi\Big)=-\frac1{2\pi}\ln L+\mathcal O(1) .
\end{equation}
At leading degree in $p$ the zero-mode operators commute and act as $a_0\to\sqrt p\,e^{\mathrm i\theta}$,
$a_0^\dagger\to\sqrt p\,e^{-\mathrm i\theta}$, so ${:}Z^2{:}\to p\,(e^{2\mathrm i\theta}+2+e^{-2\mathrm i\theta})$ and
${:}Z{:}\to\sqrt p\,(2\cos\theta)$. The closed-path amplitude is the angular average $p^3\langle(2\cos\theta)^6\rangle_\theta=
20\,p^3$. Split it by the zero-mode transfer of the inner sunset:
\begin{itemize}
\item $\Delta z_B=0$ (inner vertices $a_0$, $a_0^\dagger$ in either order, two ways) times the phase-free part of the outer
 factor $(e^{2\mathrm i\theta}+2+e^{-2\mathrm i\theta})^2$, which is $1+4+1=6$: in total 12;
\item $\Delta z_B=+2$ (inner $a_0^\dagger a_0^\dagger$) times the $e^{2\mathrm i\theta}$ coefficient of the outer factor,
 $2\cdot2=4$, and likewise $\Delta z_B=-2$: in total 8.
\end{itemize}
So $20=12+8$. The coupling factor is $g_1^4\gamma_2\gamma_3=\frac1{81}\cdot\frac12\cdot\frac23=\frac1{243}$. With
$\Pi[p^3]=-2$ the window logarithm restored in $C_4$ is
\begin{equation}
 (-2)\cdot20\cdot\Big(-\frac1{2\pi}\Big)\cdot\frac1{243}=\frac{20}{243\pi}
 =\underbrace{\frac{12}{243\pi}}_{\Delta z=0\ (=\frac4{81\pi})}+\underbrace{\frac{8}{243\pi}}_{\Delta z=\pm2} .
\end{equation}
The part that changes the zero-mode number is $\frac8{12}=\frac23$ of the diagonal part, and a diagonal window captures
$\frac{12}{20}=\frac35$ of this class. The split $20=12+8$ and the $\gamma_m$ follow from the algebra above; the factor
$\frac1{243}$ is also in \code{log\_coeffs.py}.

\begin{table}[ht]
\centering\footnotesize\renewcommand{\arraystretch}{1.15}
\begin{tabular}{@{}p{4.9cm}p{2.9cm}p{2.6cm}cccc@{}}
\toprule
class (outer state $\times$ sub-bubble) & analytic & exact fit & raw & cons & zm & $W$ \\\midrule
one $\chi$ line $\times$ vacuum bubble (also $\kappa L$) & $-\frac5{8748\pi}=-1.81933\times10^{-4}$ & $-1.81933\times10^{-4}$ & m & r & r & r \\
two $\chi$ lines $\times$ zero-mode sunset, $\Delta z=0$ & $\frac4{81\pi}=0.0157190067$ & $0.015719006698$ & m & m & r & r \\
four $\chi$ lines $\times$ one-line zero-mode bubble & $\frac5{4374\pi}=3.63866\times10^{-4}$ & $3.63865887\times10^{-4}$ & m & m & r & r \\
two $\chi$ lines $\times$ zero-mode sunset, $\Delta z=\pm2$ & $\frac8{243\pi}=0.0104793378$ & $0.010479337749$ & m & m & m & r \\
all other zero-mode classes (diagonal and off-diagonal) & 0 & $|\cdot|<2\times10^{-10}$ & & & & \\
$\chi$-attached sunsets ($D^M$) & 0 & $-2.9\times10^{-9}$; $-2.6\times10^{-7}$ & & & & \\
\bottomrule
\end{tabular}
\caption{$\ln L$ coefficients of the $C_4$ shift that each window class restores. ``m'': missing in the scheme; ``r'':
restored. The fits use the exact shifts to $L=20000$ (the $\chi$-attached class to $L=8000$) with $L_{\min}=2000$ and 4000
(\code{residual\_log/\allowbreak{}results/\allowbreak{}log\_coeffs.txt}).}
\label{htp:d2:tab:logs}
\end{table}

\paragraph{The coefficients by class.} Table~\ref{htp:d2:tab:logs} lists every class. The two other logarithmic classes follow the
same pattern. For four $\chi$ lines times a one-line bubble, the outer weight tends to $\frac14$, and when the outer state has level
$L-2j$ the inner one-line bubble keeps the budget $2j$; its window is then $-\frac1{2\pi}(\frac1j-\frac2L)$, which sums to
$-\frac1{8\pi}\ln L$; with
$g_1^4\gamma_4\gamma_1=\frac1{4374}$ and the amplitude $20p^3$ this gives $+\frac5{4374\pi}$. For one $\chi$ line times the vacuum,
Eq.~\eqref{htp:d2:eq:Ev2} gives $E_v(L)-E_v(L-2a)=-[2a(2L+6)-4a^2]/864$: the term linear in $L$ gives $\kappa L$, and the
$4a^2$ term with weight $1/(4\pi a^3)$ gives $(864\pi)^{-1}\ln L$, which the factor $-\frac{40}{81}$ of the one-line outer state (couplings,
zero-mode amplitude and $\Pi$; \code{log\_coeffs.py}) turns into $-\frac{40}{81\cdot864\pi}=-\frac5{8748\pi}$ in $C_4$. The $\ln L$ coefficient of each scheme is minus the sum of the logarithms it misses:
\begin{align}
 \beta_{\rm raw}&=-\Big(-\frac5{8748\pi}+\frac{20}{243\pi}+\frac5{4374\pi}\Big)=-\frac{725}{8748\pi},\qquad
 \beta_{\rm cons}=-\Big(\frac{20}{243\pi}+\frac5{4374\pi}\Big)=-\frac{365}{4374\pi},\nonumber\\
 \beta_{\rm zm}&=-\frac8{243\pi},\qquad \beta_W=0 ,
\end{align}
in agreement with the fitted values of Section~\ref{htp:d2:c4} to $10^{-6}$ or better. The $\chi$-attached class has no net
logarithm, but its two parts do: $+0.0369938\ln L$ directly at $\eps^4$ and $-0.0369938\ln L$ through the retuned $g_3$. This
is a local vertex renormalization, a sunset with one hard external $\chi$ line, which the tuning of $g$ absorbs. The
decomposition must be made by window class, not by diagram topology: single four-vertex topologies carry $L$, $\ln^2L$ and
$\ln L$ terms that cancel against counterterm chains (\code{decompose.py}), and the window classes are the combinations that
are separately finite in the continuum.

\paragraph{Why the tuning cannot absorb the defect.} The two tuning conditions fix two state-independent numbers, $r_n$ and
$g_n$. A shift of $r$ or $g$ changes the reference energies by $2p\,\delta r+\frac12p(p-1)\,\delta g$, which $\Pi$ removes.
The budget defect is different in kind: it is an operator proportional to the level of the spectator quanta,
${\sim}\{\hat\ell,M\}$, and through the outer amplitude $20p^3$ it is cubic in $p$. No choice of tuning conditions can
represent it.

\paragraph{Contrast with one dimension.} In $d=1$ the cutoff is on the free energy $E^0=\ell+\frac P4$
(Eq.~\eqref{htp:eq:Lambda}). Every four-particle vacuum shell then carries the same energy, so the vacuum window is of order one
per missing shell and does not decay, and $D^{\rm conn}$ repairs it (Section~\ref{htp:sec:window}). The references have
$E^0=\frac p4$, so their own windows hold vacuum shells as well, and how many depends on where the cutoff sits relative to
them (the cutoff placement of Section~\ref{htp:d1:sec:window}). By power counting the sunset density falls like
${\rm level}^{-1/2}$ ($r_2\sim\Lambda^{1/2}$), so its budget deficit ${\sim}\ell/\Lambda^{1/2}$ vanishes: this is why the
vacuum correction suffices in $d=1$ and does not in $d=2$. The connected (``folded'') $\eps^4$ term that $D^{\rm conn}$ does
not restore also decays; its contribution to the $\eps^4$ spacing of $\partial\phi$ is $-6.8\times10^{-4}$ at
$\Lambda=16.25$ and $-3.8\times10^{-4}$ at $\Lambda=64.25$ (\code{lri\_hpc\_Hamitonian}, record
\code{analysis/\allowbreak{}identity\_checks/\allowbreak{}challenge-theory/\allowbreak{}folded\_ph0.25.json}). With the
vacuum window alone the $\eps^2$ descendant spacings vanish (Table~\ref{htp:d1:tab:ident}). In $d=2$ every structure is one
power worse: the vacuum window grows like $L$ (the $\kappa L$ of the raw scheme), the sunset window is of order one, and the
coupling window is ${\sim}\ell/L$. The references have level 0, so their own windows are empty and nothing depends on where
the cutoff sits relative to them; the zero-mode dependence that $d=1$ carries through $E^0=\ell+\frac P4$ appears in $d=2$
as the separate zero-mode window.

\subsection{The complete window at \texorpdfstring{$\eps^4$}{eps4}}\label{htp:d2:app:window}
{\small\raggedright\noindent Records: \code{residual\_log/\allowbreak{}WINDOW\_SPEC.md}, \code{residual\_log/\allowbreak{}results/\allowbreak{}brute\_vs\_codes.txt},
\code{residual\_log/\allowbreak{}C4W\_by\_L.csv}, \code{residual\_log/\allowbreak{}results/\allowbreak{}\{offdiag\_L20000,\allowbreak{}nm\_L8000\}.npz},
\code{residual\_log/\allowbreak{}fits/\allowbreak{}fits\_final.txt}; code \code{residual\_log/\allowbreak{}\{nm\_window,\allowbreak{}nm\_window\_exact,\allowbreak{}brute\_full,\allowbreak{}brute\_nm,\allowbreak{}nm\_zz\_check,\allowbreak{}offdiag\_hermitian\_check,\allowbreak{}fits\_final\}.py}.\par}

\paragraph{Principle.} Section~\ref{htp:d2:c4} defines the window, $D=D^{\rm zm}+D^M$ with Eqs.~\eqref{htp:d2:eq:Dzm}
and~\eqref{htp:d2:eq:DM}, and states its principle. For a mass insertion between states of levels $\ell_s$ and $\ell_{s'}$ the
budget is the average $L-\frac12(\ell_s+\ell_{s'})$, which comes from averaging the internal orderings of the sunset: nested
and exchange orderings for a $\chi$ hop $A^\dagger A$, and the two assignments for pair creation $A^\dagger A^\dagger$ and pair
annihilation $AA$. $h^{(2)}$ contains all inner bubble types with one to four $\chi$ lines; the zero-line part and the $r_2$,
$g_2$ terms do not depend on the level and cancel. Bra-referenced and Hermitian versions of $h^{(2)}$ give the same $C_4$ to
$10^{-16}$, because the product $\langle VG\,D^{\rm zm}GV\rangle$ through which it enters is symmetric.

\paragraph{How the window enters $C_4$.} Since $D^{\rm zm}$ vanishes on level-0 states, its only contribution is
$g_1^2\langle VG\,D^{\rm zm}GV\rangle$. $D^M$ is symmetric, but it does not vanish on $|p\rangle$, because $M'$ creates two
quanta. It therefore enters already at $\eps^3$, through $2g_1\langle VGD^M\rangle$, which changes $r_3$ and $g_3$ but not $C_3$.
At $\eps^4$
\begin{align}
 \delta\mathcal E^{(4)}_p&=g_1^2\langle VGD^MGV\rangle+2g_1\big[g_1\langle VGVGD^M\rangle-g_1V_{pp}\langle VG^2D^M\rangle
 -\tfrac14\langle VG(P-p)GD^M\rangle\big],\nonumber\\
 \delta C_4^M&=\Pi\big[\delta\mathcal E^{(4)}\big]+2g_1\,\delta g_3\,\Pi\big[\langle VGV\rangle\big] ;
\end{align}
$r_3$ drops out of $C_4$, and $\Pi$ kills $\langle H_2GD^M\rangle$ and $\langle D^MGD^M\rangle$. The complete coefficient is
\begin{equation}
 C_4^W(L)=C_4^{\rm zm}(L)+\delta C_4^{\rm off}(L)+\delta C_4^M(L),
\end{equation}
where $\delta C_4^{\rm off}$ is the shift from the elements $z'\neq z$ of $D^{\rm zm}$. All three pieces are exact for every
even $L\le2000$, and the shifts are known further, to $L=20000$ (off-diagonal) and $L=8000$ ($D^M$).

\paragraph{Derivation status.} $D^{\rm zm}$ follows from the linked-cluster argument of Appendix~\ref{htp:d2:app:mech}; it is
exact at leading degree in $p$, and the differences at subleading degree vanish as $L\to\infty$. $D^M$ is derived at leading
order in $1/L$. Its slope can be replaced by the exact sum $t_M(j)$ of the sunset contributions that the cutoff removes for a
spectator of level $2j$, with both zero-mode orderings averaged; $t_M(j)\to\frac2\pi j$ per unit amplitude
(Appendix~\ref{htp:d2:app:mech}). This changes $\delta C_4^M$ by $-4.5\times10^{-4}$ at $L=100$ and $-2.4\times10^{-4}$ at $L=200$
(\code{nm\_window\_exact.py}), that is by $\mathcal O(1/L)$, with the same limit. The Newton-coefficient form of
Appendix~\ref{htp:d2:app:Wall}, another refinement, changes $C_4$ much less ($4.9\times10^{-7}$ at $L=100$), also with the same
limit; why the two refinements differ this much at finite $L$ was not analyzed.

\begin{table}[ht]
\centering\footnotesize\renewcommand{\arraystretch}{1.15}
\begin{tabular}{@{}p{9.4cm}p{2.0cm}p{3.0cm}@{}}
\toprule
check & $L$ & agreement \\\midrule
\code{offdiag\_window.py}, diagonal part, against $C_4^{\rm zm}-C_4^{\rm raw}$ of \code{c4\_exact} & $\le2000$ & $4\times10^{-11}$ \\
\code{nm\_window.py} against \code{brute\_nm.py} (exact perturbation theory with $\{\hat\ell,M'\}$ as an operator): $\delta C_4$ and $\delta g_3$ & 2, 4, 6 & $10^{-14}$ \\
\code{brute\_full.py} (complete scheme as operators; off-diagonal window from auxiliary $|z\rangle$ channels) against $C_4^{\rm zm}$, $C_4^{\rm zm}+\delta C_4^{\rm off}$ and $C_4^W$ & 2 to 12 & $\le6.3\times10^{-14}$ \\
\code{nm\_window.py} on a workstation ($L\le400$) against the symmetry cluster ($L\le2000$, 8000) & overlap & $6\times10^{-14}$ \\
\bottomrule
\end{tabular}
\caption{Checks of the window codes against brute-force exact perturbation theory in fraction arithmetic, with $D$ added as
an explicit operator (\code{residual\_log/\allowbreak{}results/\allowbreak{}brute\_vs\_codes.txt}).}
\label{htp:d2:tab:brute}
\end{table}

\paragraph{Verification at finite $L$.} Table~\ref{htp:d2:tab:brute} lists the checks. The brute-force values of the complete
scheme are $C_4^W(2)=1.966884466461$, $C_4^W(4)=2.290042077448$, $C_4^W(6)=2.546923269518$, $C_4^W(8)=2.756245037529$,
$C_4^W(10)=2.928832263364$ and $C_4^W(12)=3.073608283249$; the comparison is limited by the float64 precision of the
large-$L$ code, about $6\times10^{-14}$.

\paragraph{What does not work.}
\begin{itemize}
\item Any diagonal window (consistent, zero-mode): it misses the transitions $\Delta z\neq0$ and the $\chi$ parts by
 construction.
\item Replacing $D^{\rm zm}$ by its leading form $-\frac{g_1^2}{3\pi}\{\hat\ell,{:}Z^2{:}\}$ (\code{nm\_zz\_check.py}). The
 difference tends to $0.042838+\frac5{4374\pi}\ln L$: the inner coupling and one-line bubbles carry finite parts from hard
 outer states and the four-line $\times$ one-line logarithm, so the exact zero-mode window is required. Only the $\chi$ part
 may use the leading form.
\item A modified tuning prescription (Appendix~\ref{htp:d2:app:mech}).
\end{itemize}

\begin{table}[ht]
\centering\footnotesize\renewcommand{\arraystretch}{1.15}
\begin{tabular}{@{}p{5.6cm}lp{4.3cm}p{3.2cm}@{}}
\toprule
scheme & $C_4(2000)$ & large $L$ & limit \\\midrule
raw & 4.29126 & $-\kappa L-\frac{725}{8748\pi}\ln L$ & diverges \\
consistent ($=$ literal at this order) & 4.73178 & $-\frac{365}{4374\pi}\ln L$ & diverges \\
zero-mode (diagonal) & 4.87359 & $-\frac8{243\pi}\ln L$ & diverges \\
$+$ off-diagonal zero-mode window & 4.97161 & no $\ln L$ & $\NTwoCfourZmOff$, $0.0102545$ short \\
$+$ mass window $D^M$ ($=W$) & 4.98172 & no $\ln L$ & $5.0488682$ \\
\midrule
continuum & & & $5.0488684314$ \\
\bottomrule
\end{tabular}
\caption{$C_4$ at $L=2000$ and its large-$L$ form, window by window (\code{residual\_log/\allowbreak{}C4W\_by\_L.csv},
\code{residual\_log/\allowbreak{}fits/\allowbreak{}fits\_final.txt}).}
\label{htp:d2:tab:c4w}
\end{table}

\paragraph{Large $L$.} Table~\ref{htp:d2:tab:c4w} builds the complete window up piece by piece. The off-diagonal shift is
$\delta C_4^{\rm off}(L)=\frac8{243\pi}\ln L+0.0172932+\mathcal O(\ln^2L/L)$, from exact values to $L=20000$. It removes the
logarithm of the zero-mode scheme but leaves a finite miss of $0.0102545$. The mass-window shift converges,
$\delta C_4^M(\infty)=0.0102543394$ with a spread of $1.5\times10^{-8}$ over twelve fits of the exact values to $L=8000$, and
supplies the missing finite part: the separately extrapolated pieces give $\NTwoCfourZmOff+\NTwoCfourChi=5.0488682$. The
approach is slow, like $\ln^2L/L$: at $L=2000$, $C_4^W-C_4^{(2)}=-0.0672$ (Fig.~\ref{htp:d2:fig}c).

\paragraph{Fits of $C_4^W$.} The data are accurate to $10^{-10}$ or better, so the uncertainty is entirely the extrapolation.
\begin{itemize}
\item The towers $T_{22}$, $T_{22}+\ln^3L/L$ and $T_{23}$ without $\ln L$, with $L_{\min}\in\{200,300,400\}$ (nine fits), give
 $5.04886823$ with a half spread of $5.6\times10^{-7}$, $2.0\times10^{-7}$ below the continuum value; this is the value
 quoted in Section~\ref{htp:d2:c4}.
\item No-log fits outside this set (towers $T_{22}+\ln^3L/L$ and $T_{23}$ at $L_{\min}=600$, $T_{33}$ at $L_{\min}=400$) move
 the limit by up to $1.8\times10^{-5}$ from the continuum value; $T_{33}$ at $L_{\min}=600$ is over-fitted ($\ln^2L/L$
 coefficient $-25$, limit $2.8\times10^{-4}$ low).
\item With a free $\ln L$ term the stable towers at $L_{\min}=200$ to 300 give its coefficient between $-1.8\times10^{-6}$ and
 $+3.5\times10^{-6}$, and the limit moves by $-4.3\times10^{-5}$ to $+2.2\times10^{-5}$ (the tower $T_{33}$,
 outside this set, gives a coefficient of $-8.1\times10^{-6}$ at $L_{\min}=200$). Richer towers at larger $L_{\min}$ become unstable, for example with a free $\ln L$ coefficient up to
 $5\times10^{-4}$ and a $\ln^2L/L$ coefficient of $+16$ to $+30$; these fits have more freedom than the data constrain.
\end{itemize}
Hence the quoted uncertainty: the half spread of the nine stable fits, rounded up to $10^{-6}$. The fits outside the stable
set show that the model dependence can reach $2\times10^{-5}$, and $4\times10^{-5}$ with a free $\ln L$ term, which the
complete window excludes by construction (Table~\ref{htp:d2:tab:logs}). The fitted tail is about $-1.40\ln^2L/L$ and comes entirely from $C_4^{\rm zm}$.
Fixing it at $-5/\pi$, the value used as a prior in the fit scripts, makes the fits 500 times worse, so that prior is not
supported. In the
fits of the raw, consistent and zero-mode schemes (Section~\ref{htp:d2:c4}) the variants with the coefficient fixed only widen the quoted spread.
